
\documentclass{article}
\usepackage{colortbl}
\usepackage{makecell}
\usepackage{multirow}
\usepackage{supertabular}

\begin{document}

\newcounter{utterance}

\centering \large Interaction Transcript for game `clean\_up', experiment `1\_medium\_3obj\_ar', episode 1 with deepseek{-}v4{-}pro.
\vspace{24pt}

{ \footnotesize  \setcounter{utterance}{1}
\setlength{\tabcolsep}{0pt}
\begin{supertabular}{c@{$\;$}|p{.15\linewidth}@{}p{.15\linewidth}p{.15\linewidth}p{.15\linewidth}p{.15\linewidth}p{.15\linewidth}}
    \# & $\;$A & \multicolumn{4}{c}{Game Master} & $\;\:$B\\
    \hline

    \theutterance \stepcounter{utterance}  
    & & \multicolumn{4}{p{0.6\linewidth}}{
        \cellcolor[rgb]{0.9,0.9,0.9}{
            \makecell[{{p{\linewidth}}}]{
                \texttt{\tiny{[P1$\langle$GM]}}
                \texttt{أنا مدير لعبتك، وأنت تلعب لعبة تعاونية مع الشبكة التالية كلوحة لعب:} \\
\\ 
\texttt{```} \\
\texttt{ 1234567} \\
\texttt{╔══╤══╤═╗} \\
\texttt{║◌◌│◌◌│◌║ 1} \\
\texttt{╟──┼──┴─╢ 2} \\
\texttt{║L◌│◌◌◌◌║ 3} \\
\texttt{╟──┘◌◌◌P║ 4} \\
\texttt{║◌◌C◌◌◌◌║ 5} \\
\texttt{║◌◌┌─┬──╢ 6} \\
\texttt{║◌◌│◌│◌◌║ 7} \\
\texttt{╚══╧═╧══╝} \\
\\ 
\texttt{```} \\
\\ 
\texttt{* تعرض الحافة العلوية إحداثيات x التي تتزايد نحو اليمين، والحافة اليمنى إحداثيات y التي تتزايد نحو الأسفل.} \\
\texttt{* الكائنات التالية موضوعة عشوائياً على شبكتك: 'C', 'L', 'P'.} \\
\\ 
\\ 
\texttt{يرى اللاعب الآخر نسخة مختلفة من لوحة اللعب، حيث توضع الكائنات في مواقع عشوائية مختلفة. لا يمكنك رؤية لوحة اللاعب الآخر، ولا يمكنه رؤية لوحتك.} \\
\\ 
\\ 
\texttt{**الهدف:**} \\
\\ 
\texttt{يجب على كلا اللاعبين تحريك الكائنات على خلفية كل منهما بحيث تنتهي الكائنات المتطابقة عند الإحداثيات نفسها. عليك التواصل مع اللاعب الآخر للاتفاق على إعداد هدف مشترك.} \\
\\ 
\\ 
\texttt{**القواعد:**} \\
\\ 
\texttt{* في كل دور، يمكنك إرسال أمر واحد فقط من الأمرين التاليين:} \\
\texttt{ 1. `قل: <رسالة>`: لإرسال رسالة (كل شيء حتى فاصل السطر التالي) إلى اللاعب الآخر. سأقوم بتمريرها إلى شريكك.} \\
\texttt{2. `حرّك: <كائن>, (<X>, <Y>)`: لتحريك كائن إلى موضع جديد، حيث `<X>` هو العمود و `<Y>` هو الصف. سأخبرك بما إذا كانت حركتك ناجحة أم لا.} \\
\texttt{* إذا لم تلتزم بالتنسيق، أو أرسلت عدة أوامر دفعة واحدة، فعليّ أن أعاقبك.} \\
\texttt{* إذا تجاوز كلا اللاعبين أكثر من 8 عقوبة، فستخسران اللعبة معاً.} \\
\texttt{* من الضروري أن تتواصل مع اللاعب الآخر بشأن حالة هدفك! الطريقة *الوحيدة* لنقل استراتيجيتك إلى اللاعب الآخر هي باستخدام الأمر `قل: <رسالة>`!} \\
\\ 
\\ 
\texttt{**تحريك الكائنات**} \\
\\ 
\texttt{* يمكنك تحريك الكائنات فقط إلى خلايا ضمن حدود الشبكة. يجب أن تكون الخلية المستهدفة فارغة، أي يجب أن تحتوي فقط على الرمز '◌'.} \\
\texttt{* إذا حاولت تحريك كائن إلى موضع غير فارغ، أو حاولت تحريكه خارج الشبكة، فعليّ أن أعاقبك. ستحصل على محاولة أخرى.} \\
\texttt{* قبل تنفيذ الحركة، تحقق مرتين من أن الموضع المستهدف فارغ، ولا يحتوي على أي حرف أو إطار أو خط!} \\
\\ 
\\ 
\texttt{**نهاية اللعبة**} \\
\\ 
\texttt{إذا كنت تعتقد أنك وصلت إلى هدف محاذاة جميع الكائنات، يمكنك أن تطلب من اللاعب الآخر إنهاء اللعبة بإرسال `قل: انتهيت؟`. إذا طلب منك اللاعب الآخر إنهاء اللعبة، ورددت بـ `قل: انتهيت!`، ستنتهي اللعبة.} \\
\\ 
\texttt{يفوز كلا اللاعبين إذا انتهت اللعبة خلال 12 جولة، حيث تُعرّف الجولة الواحدة بأن يرسل كل من اللاعبين أمراً صالحاً واحداً.} \\
\\ 
\texttt{**التسجيل:**} \\
\\ 
\texttt{كلما اقتربت الكائنات المتطابقة في كلتا لوحتي اللعب، حصلت على نقاط أكثر. العقوبات تقلل نقاطك. هل يمكنك تحطيم الرقم القياسي؟} \\
\\ 
\\ 
\texttt{يرجى إرسال رسالة إلى اللاعب الآخر لبدء اللعبة!} \\
            }
        }
    }
    & \\ \\

    \theutterance \stepcounter{utterance}  
    & \multicolumn{4}{p{0.6\linewidth}}{
        \cellcolor[rgb]{0.9,0.9,0.9}{
            \makecell[{{p{\linewidth}}}]{
                \texttt{\tiny{[P1$\rangle$GM]}}
                \texttt{قل: مرحباً! دعنا نتعاون لتحريك الكائنات إلى نفس المواقع. على لوحتي، الكائنات موضوعة كالتالي: L عند (1,3)، C عند (3,5)، و P عند (7,4). ما هي مواقع الكائنات على لوحتك؟} \\
            }
        }
    }
    & & \\ \\

    \theutterance \stepcounter{utterance}  
    & & & \multicolumn{2}{p{0.3\linewidth}}{
        \cellcolor[rgb]{0.9,0.9,0.9}{
            \makecell[{{p{\linewidth}}}]{
                \texttt{\tiny{[GM$|$GM]}}
                \texttt{قل: مرحباً! دعنا نتعاون لتحريك الكائنات إلى نفس المواقع. على لوحتي، الكائنات موضوعة كالتالي: L عند (1,3)، C عند (3,5)، و P عند (7,4). ما هي مواقع الكائنات على لوحتك؟} \\
            }
        }
    }
    & & \\ \\

    \theutterance \stepcounter{utterance}  
    & & \multicolumn{4}{p{0.6\linewidth}}{
        \cellcolor[rgb]{0.9,0.9,0.9}{
            \makecell[{{p{\linewidth}}}]{
                \texttt{\tiny{[GM$\rangle$P2]}}
                \texttt{أنا مدير لعبتك، وأنت تلعب لعبة تعاونية مع الشبكة التالية كلوحة لعب:} \\
\\ 
\texttt{```} \\
\texttt{ 1234567} \\
\texttt{╔══╤══╤═╗} \\
\texttt{║◌P│◌◌│◌║ 1} \\
\texttt{╟──┼──┴─╢ 2} \\
\texttt{║◌◌│◌L◌◌║ 3} \\
\texttt{╟──┘◌◌◌◌║ 4} \\
\texttt{║◌◌◌◌◌◌C║ 5} \\
\texttt{║◌◌┌─┬──╢ 6} \\
\texttt{║◌◌│◌│◌◌║ 7} \\
\texttt{╚══╧═╧══╝} \\
\\ 
\texttt{```} \\
\\ 
\texttt{* تعرض الحافة العلوية إحداثيات x التي تتزايد نحو اليمين، والحافة اليمنى إحداثيات y التي تتزايد نحو الأسفل.} \\
\texttt{* الكائنات التالية موضوعة عشوائياً على شبكتك: 'C', 'L', 'P'.} \\
\\ 
\\ 
\texttt{يرى اللاعب الآخر نسخة مختلفة من لوحة اللعب، حيث توضع الكائنات في مواقع عشوائية مختلفة. لا يمكنك رؤية لوحة اللاعب الآخر، ولا يمكنه رؤية لوحتك.} \\
\\ 
\\ 
\texttt{**الهدف:**} \\
\\ 
\texttt{يجب على كلا اللاعبين تحريك الكائنات على خلفية كل منهما بحيث تنتهي الكائنات المتطابقة عند الإحداثيات نفسها. عليك التواصل مع اللاعب الآخر للاتفاق على إعداد هدف مشترك.} \\
\\ 
\\ 
\texttt{**القواعد:**} \\
\\ 
\texttt{* في كل دور، يمكنك إرسال أمر واحد فقط من الأمرين التاليين:} \\
\texttt{ 1. `قل: <رسالة>`: لإرسال رسالة (كل شيء حتى فاصل السطر التالي) إلى اللاعب الآخر. سأقوم بتمريرها إلى شريكك.} \\
\texttt{2. `حرّك: <كائن>, (<X>, <Y>)`: لتحريك كائن إلى موضع جديد، حيث `<X>` هو العمود و `<Y>` هو الصف. سأخبرك بما إذا كانت حركتك ناجحة أم لا.} \\
\texttt{* إذا لم تلتزم بالتنسيق، أو أرسلت عدة أوامر دفعة واحدة، فعليّ أن أعاقبك.} \\
\texttt{* إذا تجاوز كلا اللاعبين أكثر من 8 عقوبة، فستخسران اللعبة معاً.} \\
\texttt{* من الضروري أن تتواصل مع اللاعب الآخر بشأن حالة هدفك! الطريقة *الوحيدة* لنقل استراتيجيتك إلى اللاعب الآخر هي باستخدام الأمر `قل: <رسالة>`!} \\
\\ 
\\ 
\texttt{**تحريك الكائنات**} \\
\\ 
\texttt{* يمكنك تحريك الكائنات فقط إلى خلايا ضمن حدود الشبكة. يجب أن تكون الخلية المستهدفة فارغة، أي يجب أن تحتوي فقط على الرمز '◌'.} \\
\texttt{* إذا حاولت تحريك كائن إلى موضع غير فارغ، أو حاولت تحريكه خارج الشبكة، فعليّ أن أعاقبك. ستحصل على محاولة أخرى.} \\
\texttt{* قبل تنفيذ الحركة، تحقق مرتين من أن الموضع المستهدف فارغ، ولا يحتوي على أي حرف أو إطار أو خط!} \\
\\ 
\\ 
\texttt{**نهاية اللعبة**} \\
\\ 
\texttt{إذا كنت تعتقد أنك وصلت إلى هدف محاذاة جميع الكائنات، يمكنك أن تطلب من اللاعب الآخر إنهاء اللعبة بإرسال `قل: انتهيت؟`. إذا طلب منك اللاعب الآخر إنهاء اللعبة، ورددت بـ `قل: انتهيت!`، ستنتهي اللعبة.} \\
\\ 
\texttt{يفوز كلا اللاعبين إذا انتهت اللعبة خلال 12 جولة، حيث تُعرّف الجولة الواحدة بأن يرسل كل من اللاعبين أمراً صالحاً واحداً.} \\
\\ 
\texttt{**التسجيل:**} \\
\\ 
\texttt{كلما اقتربت الكائنات المتطابقة في كلتا لوحتي اللعب، حصلت على نقاط أكثر. العقوبات تقلل نقاطك. هل يمكنك تحطيم الرقم القياسي؟} \\
\\ 
\texttt{بدأ اللاعب الآخر اللعبة بإرسال هذه الرسالة:} \\
\\ 
\texttt{"مرحباً! دعنا نتعاون لتحريك الكائنات إلى نفس المواقع. على لوحتي، الكائنات موضوعة كالتالي: L عند (1,3)، C عند (3,5)، و P عند (7,4). ما هي مواقع الكائنات على لوحتك؟"} \\
\\ 
\texttt{ما هو أمرك الأول؟} \\
            }
        }
    }
    & \\ \\

    \theutterance \stepcounter{utterance}  
    & & & \multicolumn{4}{p{0.6\linewidth}}{
        \cellcolor[rgb]{0.9,0.9,0.9}{
            \makecell[{{p{\linewidth}}}]{
                \texttt{\tiny{[GM$\langle$P2]}}
                \texttt{قل: مرحباً! على لوحتي، الكائنات موضوعة كالتالي: P عند (2,1)، L عند (4,3)، و C عند (7,5). لنعمل معاً لإيجاد مواقع مشتركة.} \\
            }
        }
    }
     \\ \\

    \theutterance \stepcounter{utterance}  
    & & & \multicolumn{2}{p{0.3\linewidth}}{
        \cellcolor[rgb]{0.9,0.9,0.9}{
            \makecell[{{p{\linewidth}}}]{
                \texttt{\tiny{[GM$|$GM]}}
                \texttt{قل: مرحباً! على لوحتي، الكائنات موضوعة كالتالي: P عند (2,1)، L عند (4,3)، و C عند (7,5). لنعمل معاً لإيجاد مواقع مشتركة.} \\
            }
        }
    }
    & & \\ \\

    \theutterance \stepcounter{utterance}  
    & & \multicolumn{4}{p{0.6\linewidth}}{
        \cellcolor[rgb]{0.9,0.9,0.9}{
            \makecell[{{p{\linewidth}}}]{
                \texttt{\tiny{[P1$\langle$GM]}}
                \texttt{تم تمرير رسالتك إلى اللاعب الآخر.} \\
\\ 
\\ 
\texttt{أنت تلعب حالياً الجولة 1 من أصل 12 كحد أقصى.} \\
\texttt{لقد جمعتما معاً 0 من أصل 8 عقوبة.} \\
\texttt{أرسل اللاعب الآخر هذه الرسالة:} \\
\\ 
\\ 
\texttt{"مرحباً! على لوحتي، الكائنات موضوعة كالتالي: P عند (2,1)، L عند (4,3)، و C عند (7,5). لنعمل معاً لإيجاد مواقع مشتركة."} \\
\\ 
\\ 
\texttt{ما هو أمرك التالي؟} \\
            }
        }
    }
    & \\ \\

    \theutterance \stepcounter{utterance}  
    & \multicolumn{4}{p{0.6\linewidth}}{
        \cellcolor[rgb]{0.9,0.9,0.9}{
            \makecell[{{p{\linewidth}}}]{
                \texttt{\tiny{[P1$\rangle$GM]}}
                \texttt{قل: حسناً، لنقارن المواقع. لدينا ثلاثة كائنات: C، L، و P. على لوحتي: L(1,3) و C(3,5) و P(7,4). على لوحتك: P(2,1) و L(4,3) و C(7,5). نحتاج لإيجاد ثلاثة مواقع فارغة على كلا اللوحتين لننقل إليها الكائنات. هل يمكنك التحقق من الخلايا الفارغة التي تناسبك؟ أقترح أن نبدأ بتحديد مواقع مستهدفة فارغة في كلا اللوحتين مثل (2,2) و (5,3) و (6,6). ما رأيك؟} \\
            }
        }
    }
    & & \\ \\

    \theutterance \stepcounter{utterance}  
    & & & \multicolumn{2}{p{0.3\linewidth}}{
        \cellcolor[rgb]{0.9,0.9,0.9}{
            \makecell[{{p{\linewidth}}}]{
                \texttt{\tiny{[GM$|$GM]}}
                \texttt{قل: حسناً، لنقارن المواقع. لدينا ثلاثة كائنات: C، L، و P. على لوحتي: L(1,3) و C(3,5) و P(7,4). على لوحتك: P(2,1) و L(4,3) و C(7,5). نحتاج لإيجاد ثلاثة مواقع فارغة على كلا اللوحتين لننقل إليها الكائنات. هل يمكنك التحقق من الخلايا الفارغة التي تناسبك؟ أقترح أن نبدأ بتحديد مواقع مستهدفة فارغة في كلا اللوحتين مثل (2,2) و (5,3) و (6,6). ما رأيك؟} \\
            }
        }
    }
    & & \\ \\

    \theutterance \stepcounter{utterance}  
    & & \multicolumn{4}{p{0.6\linewidth}}{
        \cellcolor[rgb]{0.9,0.9,0.9}{
            \makecell[{{p{\linewidth}}}]{
                \texttt{\tiny{[GM$\rangle$P2]}}
                \texttt{تم تمرير رسالتك إلى اللاعب الآخر.} \\
\\ 
\\ 
\texttt{أنت تلعب حالياً الجولة 2 من أصل 12 كحد أقصى.} \\
\texttt{لقد جمعتما معاً 0 من أصل 8 عقوبة.} \\
\texttt{أرسل اللاعب الآخر هذه الرسالة:} \\
\\ 
\\ 
\texttt{"حسناً، لنقارن المواقع. لدينا ثلاثة كائنات: C، L، و P. على لوحتي: L(1,3) و C(3,5) و P(7,4). على لوحتك: P(2,1) و L(4,3) و C(7,5). نحتاج لإيجاد ثلاثة مواقع فارغة على كلا اللوحتين لننقل إليها الكائنات. هل يمكنك التحقق من الخلايا الفارغة التي تناسبك؟ أقترح أن نبدأ بتحديد مواقع مستهدفة فارغة في كلا اللوحتين مثل (2,2) و (5,3) و (6,6). ما رأيك؟"} \\
\\ 
\\ 
\texttt{ما هو أمرك التالي؟} \\
            }
        }
    }
    & \\ \\

    \theutterance \stepcounter{utterance}  
    & & & \multicolumn{4}{p{0.6\linewidth}}{
        \cellcolor[rgb]{0.9,0.9,0.9}{
            \makecell[{{p{\linewidth}}}]{
                \texttt{\tiny{[GM$\langle$P2]}}
                \texttt{قل: فكرة جيدة! دعني أتحقق من الخلايا التي اقترحتها على لوحتي. (2,2) فارغة، (5,3) فارغة، و(6,6) فارغة أيضًا. جميعها متاحة للنقل. هل هذه الخلايا فارغة أيضًا على لوحتك؟ إذا كان الأمر كذلك، فلنحدد أي كائن يذهب إلى أي موقع.} \\
            }
        }
    }
     \\ \\

    \theutterance \stepcounter{utterance}  
    & & & \multicolumn{2}{p{0.3\linewidth}}{
        \cellcolor[rgb]{0.9,0.9,0.9}{
            \makecell[{{p{\linewidth}}}]{
                \texttt{\tiny{[GM$|$GM]}}
                \texttt{قل: فكرة جيدة! دعني أتحقق من الخلايا التي اقترحتها على لوحتي. (2,2) فارغة، (5,3) فارغة، و(6,6) فارغة أيضًا. جميعها متاحة للنقل. هل هذه الخلايا فارغة أيضًا على لوحتك؟ إذا كان الأمر كذلك، فلنحدد أي كائن يذهب إلى أي موقع.} \\
            }
        }
    }
    & & \\ \\

    \theutterance \stepcounter{utterance}  
    & & \multicolumn{4}{p{0.6\linewidth}}{
        \cellcolor[rgb]{0.9,0.9,0.9}{
            \makecell[{{p{\linewidth}}}]{
                \texttt{\tiny{[P1$\langle$GM]}}
                \texttt{تم تمرير رسالتك إلى اللاعب الآخر.} \\
\\ 
\\ 
\texttt{أنت تلعب حالياً الجولة 2 من أصل 12 كحد أقصى.} \\
\texttt{لقد جمعتما معاً 0 من أصل 8 عقوبة.} \\
\texttt{أرسل اللاعب الآخر هذه الرسالة:} \\
\\ 
\\ 
\texttt{"فكرة جيدة! دعني أتحقق من الخلايا التي اقترحتها على لوحتي. (2,2) فارغة، (5,3) فارغة، و(6,6) فارغة أيضًا. جميعها متاحة للنقل. هل هذه الخلايا فارغة أيضًا على لوحتك؟ إذا كان الأمر كذلك، فلنحدد أي كائن يذهب إلى أي موقع."} \\
\\ 
\\ 
\texttt{ما هو أمرك التالي؟} \\
            }
        }
    }
    & \\ \\

    \theutterance \stepcounter{utterance}  
    & \multicolumn{4}{p{0.6\linewidth}}{
        \cellcolor[rgb]{0.9,0.9,0.9}{
            \makecell[{{p{\linewidth}}}]{
                \texttt{\tiny{[P1$\rangle$GM]}}
                \texttt{لنتحقق من الخلايا على لوحتي: (2,2) تناظر الموضع في الصف 2 والعمود 2، وهو ضمن المنطقة الفارغة في الأعلى {-} نعم، إنها فارغة. (5,3): الصف 3 والعمود 5 {-} أرى أن الصف 3 به L عند (1,3) ثم فراغات حتى النهاية، لذا (5,3) فارغة. (6,6): الصف 6 والعمود 6 {-} الصف 6 به خلايا فارغة بعد العمود 4، لذا (6,6) فارغة أيضاً. جميعها متاحة!} \\
\\ 
\texttt{الآن لنوزع الكائنات. أقترح:} \\
\texttt{{-} C إلى (2,2)} \\
\texttt{{-} L إلى (5,3)} \\
\texttt{{-} P إلى (6,6)} \\
\\ 
\texttt{هذا يعني تحريك C من (3,5) إلى (2,2)، و L من (1,3) إلى (5,3)، و P من (7,4) إلى (6,6). هل توافق على هذه التوزيعة؟ إذا وافقت، يمكننا البدء بالتحريك.} \\
            }
        }
    }
    & & \\ \\

    \theutterance \stepcounter{utterance}  
    & & & \multicolumn{2}{p{0.3\linewidth}}{
        \cellcolor[rgb]{0.9,0.9,0.9}{
            \makecell[{{p{\linewidth}}}]{
                \texttt{\tiny{[GM$|$GM]}}
                \texttt{لنتحقق من الخلايا على لوحتي: (2,2) تناظر الموضع في الصف 2 والعمود 2، وهو ضمن المنطقة الفارغة في الأعلى {-} نعم، إنها فارغة. (5,3): الصف 3 والعمود 5 {-} أرى أن الصف 3 به L عند (1,3) ثم فراغات حتى النهاية، لذا (5,3) فارغة. (6,6): الصف 6 والعمود 6 {-} الصف 6 به خلايا فارغة بعد العمود 4، لذا (6,6) فارغة أيضاً. جميعها متاحة!} \\
\\ 
\texttt{الآن لنوزع الكائنات. أقترح:} \\
\texttt{{-} C إلى (2,2)} \\
\texttt{{-} L إلى (5,3)} \\
\texttt{{-} P إلى (6,6)} \\
\\ 
\texttt{هذا يعني تحريك C من (3,5) إلى (2,2)، و L من (1,3) إلى (5,3)، و P من (7,4) إلى (6,6). هل توافق على هذه التوزيعة؟ إذا وافقت، يمكننا البدء بالتحريك.} \\
            }
        }
    }
    & & \\ \\

    \theutterance \stepcounter{utterance}  
    & & & \multicolumn{2}{p{0.3\linewidth}}{
        \cellcolor[rgb]{0.9,0.9,0.9}{
            \makecell[{{p{\linewidth}}}]{
                \texttt{\tiny{[GM$|$GM]}}
                \texttt{Invalid response format} \\
            }
        }
    }
    & & \\ \\

    \theutterance \stepcounter{utterance}  
    & & \multicolumn{4}{p{0.6\linewidth}}{
        \cellcolor[rgb]{0.9,0.9,0.9}{
            \makecell[{{p{\linewidth}}}]{
                \texttt{\tiny{[P1$\langle$GM]}}
                \texttt{عقوبة: رسالتك ليست بالتنسيق المتوقع!} \\
\\ 
\texttt{تأكد من أن ردك يحتوي فقط على `قل: <رسالة>` أو `حرّك: <كائن>, (<X>, <Y>)`، ولا شيء غير ذلك!} \\
\\ 
\texttt{لقد جمعتما معاً 1 من أصل 8 عقوبة.} \\
\texttt{يرجى المحاولة مرة أخرى!} \\
            }
        }
    }
    & \\ \\

    \theutterance \stepcounter{utterance}  
    & \multicolumn{4}{p{0.6\linewidth}}{
        \cellcolor[rgb]{0.9,0.9,0.9}{
            \makecell[{{p{\linewidth}}}]{
                \texttt{\tiny{[P1$\rangle$GM]}}
                \texttt{قل: عظيم! جميع الخلايا (2,2) و (5,3) و (6,6) فارغة على لوحتي أيضاً. التوزيع المقترح: C إلى (2,2)، L إلى (5,3)، P إلى (6,6). هل توافق؟ إذا وافقت، سنبدأ بالتحريك.} \\
            }
        }
    }
    & & \\ \\

    \theutterance \stepcounter{utterance}  
    & & & \multicolumn{2}{p{0.3\linewidth}}{
        \cellcolor[rgb]{0.9,0.9,0.9}{
            \makecell[{{p{\linewidth}}}]{
                \texttt{\tiny{[GM$|$GM]}}
                \texttt{قل: عظيم! جميع الخلايا (2,2) و (5,3) و (6,6) فارغة على لوحتي أيضاً. التوزيع المقترح: C إلى (2,2)، L إلى (5,3)، P إلى (6,6). هل توافق؟ إذا وافقت، سنبدأ بالتحريك.} \\
            }
        }
    }
    & & \\ \\

    \theutterance \stepcounter{utterance}  
    & & \multicolumn{4}{p{0.6\linewidth}}{
        \cellcolor[rgb]{0.9,0.9,0.9}{
            \makecell[{{p{\linewidth}}}]{
                \texttt{\tiny{[GM$\rangle$P2]}}
                \texttt{تم تمرير رسالتك إلى اللاعب الآخر.} \\
\\ 
\\ 
\texttt{أنت تلعب حالياً الجولة 3 من أصل 12 كحد أقصى.} \\
\texttt{لقد جمعتما معاً 1 من أصل 8 عقوبة.} \\
\texttt{أرسل اللاعب الآخر هذه الرسالة:} \\
\\ 
\\ 
\texttt{"عظيم! جميع الخلايا (2,2) و (5,3) و (6,6) فارغة على لوحتي أيضاً. التوزيع المقترح: C إلى (2,2)، L إلى (5,3)، P إلى (6,6). هل توافق؟ إذا وافقت، سنبدأ بالتحريك."} \\
\\ 
\\ 
\texttt{ما هو أمرك التالي؟} \\
            }
        }
    }
    & \\ \\

    \theutterance \stepcounter{utterance}  
    & & & \multicolumn{4}{p{0.6\linewidth}}{
        \cellcolor[rgb]{0.9,0.9,0.9}{
            \makecell[{{p{\linewidth}}}]{
                \texttt{\tiny{[GM$\langle$P2]}}
                \texttt{حرّك: C, (2,2)} \\
            }
        }
    }
     \\ \\

    \theutterance \stepcounter{utterance}  
    & & & \multicolumn{2}{p{0.3\linewidth}}{
        \cellcolor[rgb]{0.9,0.9,0.9}{
            \makecell[{{p{\linewidth}}}]{
                \texttt{\tiny{[GM$|$GM]}}
                \texttt{حرّك: C, (2,2)} \\
            }
        }
    }
    & & \\ \\

    \theutterance \stepcounter{utterance}  
    & & & \multicolumn{2}{p{0.3\linewidth}}{
        \cellcolor[rgb]{0.9,0.9,0.9}{
            \makecell[{{p{\linewidth}}}]{
                \texttt{\tiny{[GM$|$GM]}}
                \texttt{عقوبة: (2,2) ليست فارغة، بل تحتوي على '─'.} \\
\\ 
\texttt{لقد جمعتما معاً 2 من أصل 8 عقوبة.} \\
\texttt{يرجى المحاولة مرة أخرى!} \\
            }
        }
    }
    & & \\ \\

    \theutterance \stepcounter{utterance}  
    & & \multicolumn{4}{p{0.6\linewidth}}{
        \cellcolor[rgb]{0.9,0.9,0.9}{
            \makecell[{{p{\linewidth}}}]{
                \texttt{\tiny{[GM$\rangle$P2]}}
                \texttt{عقوبة: (2,2) ليست فارغة، بل تحتوي على '─'.} \\
\\ 
\texttt{لقد جمعتما معاً 2 من أصل 8 عقوبة.} \\
\texttt{يرجى المحاولة مرة أخرى!} \\
            }
        }
    }
    & \\ \\

    \theutterance \stepcounter{utterance}  
    & & & \multicolumn{4}{p{0.6\linewidth}}{
        \cellcolor[rgb]{0.9,0.9,0.9}{
            \makecell[{{p{\linewidth}}}]{
                \texttt{\tiny{[GM$\langle$P2]}}
                \texttt{عذراً على الخطأ! دعني أتحقق من لوحتي مرة أخرى. أرى الآن أن (2,2) تحتوي على '─' بالفعل. الخلايا الفارغة المتاحة لدي هي: (2,4) و(2,6) و(3,2) و(3,4) و(3,6) و(4,2) و(4,6) و(5,2) و(5,4) و(6,2) و(6,4) و(6,5). هل يمكنك التحقق من هذه الخلايا على لوحتك لنتمكن من اختيار ثلاثة منها معاً؟} \\
            }
        }
    }
     \\ \\

    \theutterance \stepcounter{utterance}  
    & & & \multicolumn{2}{p{0.3\linewidth}}{
        \cellcolor[rgb]{0.9,0.9,0.9}{
            \makecell[{{p{\linewidth}}}]{
                \texttt{\tiny{[GM$|$GM]}}
                \texttt{عذراً على الخطأ! دعني أتحقق من لوحتي مرة أخرى. أرى الآن أن (2,2) تحتوي على '─' بالفعل. الخلايا الفارغة المتاحة لدي هي: (2,4) و(2,6) و(3,2) و(3,4) و(3,6) و(4,2) و(4,6) و(5,2) و(5,4) و(6,2) و(6,4) و(6,5). هل يمكنك التحقق من هذه الخلايا على لوحتك لنتمكن من اختيار ثلاثة منها معاً؟} \\
            }
        }
    }
    & & \\ \\

    \theutterance \stepcounter{utterance}  
    & & & \multicolumn{2}{p{0.3\linewidth}}{
        \cellcolor[rgb]{0.9,0.9,0.9}{
            \makecell[{{p{\linewidth}}}]{
                \texttt{\tiny{[GM$|$GM]}}
                \texttt{Invalid response format} \\
            }
        }
    }
    & & \\ \\

    \theutterance \stepcounter{utterance}  
    & & \multicolumn{4}{p{0.6\linewidth}}{
        \cellcolor[rgb]{0.9,0.9,0.9}{
            \makecell[{{p{\linewidth}}}]{
                \texttt{\tiny{[GM$\rangle$P2]}}
                \texttt{عقوبة: رسالتك ليست بالتنسيق المتوقع!} \\
\\ 
\texttt{تأكد من أن ردك يحتوي فقط على `قل: <رسالة>` أو `حرّك: <كائن>, (<X>, <Y>)`، ولا شيء غير ذلك!} \\
\\ 
\texttt{لقد جمعتما معاً 3 من أصل 8 عقوبة.} \\
\texttt{يرجى المحاولة مرة أخرى!} \\
            }
        }
    }
    & \\ \\

    \theutterance \stepcounter{utterance}  
    & & & \multicolumn{4}{p{0.6\linewidth}}{
        \cellcolor[rgb]{0.9,0.9,0.9}{
            \makecell[{{p{\linewidth}}}]{
                \texttt{\tiny{[GM$\langle$P2]}}
                \texttt{قل: عذراً على الخطأ! دعني أتحقق من لوحتي مرة أخرى. أرى الآن أن (2,2) تحتوي على '─' بالفعل. الخلايا الفارغة المتاحة لدي هي: (2,4) و(2,6) و(3,2) و(3,4) و(3,6) و(4,2) و(4,6) و(5,2) و(5,4) و(6,2) و(6,4) و(6,5). هل يمكنك التحقق من هذه الخلايا على لوحتك لنتمكن من اختيار ثلاثة منها معاً؟} \\
            }
        }
    }
     \\ \\

    \theutterance \stepcounter{utterance}  
    & & & \multicolumn{2}{p{0.3\linewidth}}{
        \cellcolor[rgb]{0.9,0.9,0.9}{
            \makecell[{{p{\linewidth}}}]{
                \texttt{\tiny{[GM$|$GM]}}
                \texttt{قل: عذراً على الخطأ! دعني أتحقق من لوحتي مرة أخرى. أرى الآن أن (2,2) تحتوي على '─' بالفعل. الخلايا الفارغة المتاحة لدي هي: (2,4) و(2,6) و(3,2) و(3,4) و(3,6) و(4,2) و(4,6) و(5,2) و(5,4) و(6,2) و(6,4) و(6,5). هل يمكنك التحقق من هذه الخلايا على لوحتك لنتمكن من اختيار ثلاثة منها معاً؟} \\
            }
        }
    }
    & & \\ \\

    \theutterance \stepcounter{utterance}  
    & & \multicolumn{4}{p{0.6\linewidth}}{
        \cellcolor[rgb]{0.9,0.9,0.9}{
            \makecell[{{p{\linewidth}}}]{
                \texttt{\tiny{[P1$\langle$GM]}}
                \texttt{تم تمرير رسالتك إلى اللاعب الآخر.} \\
\\ 
\\ 
\texttt{أنت تلعب حالياً الجولة 3 من أصل 12 كحد أقصى.} \\
\texttt{لقد جمعتما معاً 3 من أصل 8 عقوبة.} \\
\texttt{أرسل اللاعب الآخر هذه الرسالة:} \\
\\ 
\\ 
\texttt{"عذراً على الخطأ! دعني أتحقق من لوحتي مرة أخرى. أرى الآن أن (2,2) تحتوي على '─' بالفعل. الخلايا الفارغة المتاحة لدي هي: (2,4) و(2,6) و(3,2) و(3,4) و(3,6) و(4,2) و(4,6) و(5,2) و(5,4) و(6,2) و(6,4) و(6,5). هل يمكنك التحقق من هذه الخلايا على لوحتك لنتمكن من اختيار ثلاثة منها معاً؟"} \\
\\ 
\\ 
\texttt{ما هو أمرك التالي؟} \\
            }
        }
    }
    & \\ \\

    \theutterance \stepcounter{utterance}  
    & \multicolumn{4}{p{0.6\linewidth}}{
        \cellcolor[rgb]{0.9,0.9,0.9}{
            \makecell[{{p{\linewidth}}}]{
                \texttt{\tiny{[P1$\rangle$GM]}}
                \texttt{قل: لا تقلق! دعني أتحقق من الخلايا التي ذكرتها على لوحتي. سأفحص كل واحدة:} \\
\texttt{(2,4): فارغة} \\
\texttt{(2,6): فارغة} \\
\texttt{(3,2): فارغة} \\
\texttt{(3,4): فارغة} \\
\texttt{(3,6): فارغة} \\
\texttt{(4,2): فارغة} \\
\texttt{(4,6): فارغة} \\
\texttt{(5,2): فارغة} \\
\texttt{(5,4): فارغة} \\
\texttt{(6,2): فارغة} \\
\texttt{(6,4): فارغة} \\
\texttt{(6,5): فارغة} \\
\\ 
\texttt{جميع الخلايا التي ذكرتها فارغة على لوحتي أيضاً! فلنختر ثلاثة منها. أقترح: C إلى (3,2)، L إلى (5,4)، P إلى (6,5). ما رأيك؟} \\
            }
        }
    }
    & & \\ \\

    \theutterance \stepcounter{utterance}  
    & & & \multicolumn{2}{p{0.3\linewidth}}{
        \cellcolor[rgb]{0.9,0.9,0.9}{
            \makecell[{{p{\linewidth}}}]{
                \texttt{\tiny{[GM$|$GM]}}
                \texttt{قل: لا تقلق! دعني أتحقق من الخلايا التي ذكرتها على لوحتي. سأفحص كل واحدة:} \\
\texttt{(2,4): فارغة} \\
\texttt{(2,6): فارغة} \\
\texttt{(3,2): فارغة} \\
\texttt{(3,4): فارغة} \\
\texttt{(3,6): فارغة} \\
\texttt{(4,2): فارغة} \\
\texttt{(4,6): فارغة} \\
\texttt{(5,2): فارغة} \\
\texttt{(5,4): فارغة} \\
\texttt{(6,2): فارغة} \\
\texttt{(6,4): فارغة} \\
\texttt{(6,5): فارغة} \\
\\ 
\texttt{جميع الخلايا التي ذكرتها فارغة على لوحتي أيضاً! فلنختر ثلاثة منها. أقترح: C إلى (3,2)، L إلى (5,4)، P إلى (6,5). ما رأيك؟} \\
            }
        }
    }
    & & \\ \\

    \theutterance \stepcounter{utterance}  
    & & & \multicolumn{2}{p{0.3\linewidth}}{
        \cellcolor[rgb]{0.9,0.9,0.9}{
            \makecell[{{p{\linewidth}}}]{
                \texttt{\tiny{[GM$|$GM]}}
                \texttt{Invalid format: tail is not empty: '} \\
\texttt{(2,4): فارغة} \\
\texttt{(2,6): فارغة} \\
\texttt{(3,2): فارغة} \\
\texttt{(3,4): فارغة} \\
\texttt{(3,6): فارغة} \\
\texttt{(4,2): فارغة} \\
\texttt{(4,6): فارغة} \\
\texttt{(5,2): فارغة} \\
\texttt{(5,4): فارغة} \\
\texttt{(6,2): فارغة} \\
\texttt{(6,4): فارغة} \\
\texttt{(6,5): فارغة} \\
\\ 
\texttt{جميع الخلايا التي ذكرتها فارغة على لوحتي أيضاً! فلنختر ثلاثة منها. أقترح: C إلى (3,2)، L إلى (5,4)، P إلى (6,5). ما رأيك؟} \\
            }
        }
    }
    & & \\ \\

    \theutterance \stepcounter{utterance}  
    & & \multicolumn{4}{p{0.6\linewidth}}{
        \cellcolor[rgb]{0.9,0.9,0.9}{
            \makecell[{{p{\linewidth}}}]{
                \texttt{\tiny{[P1$\langle$GM]}}
                \texttt{عقوبة: يجب ألا تحتوي رسالتك على أي شيء بعد الأمر!} \\
\\ 
\texttt{تأكد من أن ردك يحتوي فقط على `قل: <رسالة>` أو `حرّك: <كائن>, (<X>, <Y>)`، ولا شيء غير ذلك!} \\
\\ 
\texttt{لقد جمعتما معاً 4 من أصل 8 عقوبة.} \\
\texttt{يرجى المحاولة مرة أخرى!} \\
            }
        }
    }
    & \\ \\

    \theutterance \stepcounter{utterance}  
    & \multicolumn{4}{p{0.6\linewidth}}{
        \cellcolor[rgb]{0.9,0.9,0.9}{
            \makecell[{{p{\linewidth}}}]{
                \texttt{\tiny{[P1$\rangle$GM]}}
                \texttt{قل: جميع الخلايا التي ذكرتها فارغة على لوحتي. فلنختر: C إلى (3,2)، L إلى (5,4)، P إلى (6,5). هل توافق؟} \\
            }
        }
    }
    & & \\ \\

    \theutterance \stepcounter{utterance}  
    & & & \multicolumn{2}{p{0.3\linewidth}}{
        \cellcolor[rgb]{0.9,0.9,0.9}{
            \makecell[{{p{\linewidth}}}]{
                \texttt{\tiny{[GM$|$GM]}}
                \texttt{قل: جميع الخلايا التي ذكرتها فارغة على لوحتي. فلنختر: C إلى (3,2)، L إلى (5,4)، P إلى (6,5). هل توافق؟} \\
            }
        }
    }
    & & \\ \\

    \theutterance \stepcounter{utterance}  
    & & \multicolumn{4}{p{0.6\linewidth}}{
        \cellcolor[rgb]{0.9,0.9,0.9}{
            \makecell[{{p{\linewidth}}}]{
                \texttt{\tiny{[GM$\rangle$P2]}}
                \texttt{تم تمرير رسالتك إلى اللاعب الآخر.} \\
\\ 
\\ 
\texttt{أنت تلعب حالياً الجولة 4 من أصل 12 كحد أقصى.} \\
\texttt{لقد جمعتما معاً 4 من أصل 8 عقوبة.} \\
\texttt{أرسل اللاعب الآخر هذه الرسالة:} \\
\\ 
\\ 
\texttt{"جميع الخلايا التي ذكرتها فارغة على لوحتي. فلنختر: C إلى (3,2)، L إلى (5,4)، P إلى (6,5). هل توافق؟"} \\
\\ 
\\ 
\texttt{ما هو أمرك التالي؟} \\
            }
        }
    }
    & \\ \\

    \theutterance \stepcounter{utterance}  
    & & & \multicolumn{4}{p{0.6\linewidth}}{
        \cellcolor[rgb]{0.9,0.9,0.9}{
            \makecell[{{p{\linewidth}}}]{
                \texttt{\tiny{[GM$\langle$P2]}}
                \texttt{حرّك: C, (3,2)} \\
            }
        }
    }
     \\ \\

    \theutterance \stepcounter{utterance}  
    & & & \multicolumn{2}{p{0.3\linewidth}}{
        \cellcolor[rgb]{0.9,0.9,0.9}{
            \makecell[{{p{\linewidth}}}]{
                \texttt{\tiny{[GM$|$GM]}}
                \texttt{حرّك: C, (3,2)} \\
            }
        }
    }
    & & \\ \\

    \theutterance \stepcounter{utterance}  
    & & & \multicolumn{2}{p{0.3\linewidth}}{
        \cellcolor[rgb]{0.9,0.9,0.9}{
            \makecell[{{p{\linewidth}}}]{
                \texttt{\tiny{[GM$|$GM]}}
                \texttt{عقوبة: (3,2) ليست فارغة، بل تحتوي على '┼'.} \\
\\ 
\texttt{لقد جمعتما معاً 5 من أصل 8 عقوبة.} \\
\texttt{يرجى المحاولة مرة أخرى!} \\
            }
        }
    }
    & & \\ \\

    \theutterance \stepcounter{utterance}  
    & & \multicolumn{4}{p{0.6\linewidth}}{
        \cellcolor[rgb]{0.9,0.9,0.9}{
            \makecell[{{p{\linewidth}}}]{
                \texttt{\tiny{[GM$\rangle$P2]}}
                \texttt{عقوبة: (3,2) ليست فارغة، بل تحتوي على '┼'.} \\
\\ 
\texttt{لقد جمعتما معاً 5 من أصل 8 عقوبة.} \\
\texttt{يرجى المحاولة مرة أخرى!} \\
            }
        }
    }
    & \\ \\

    \theutterance \stepcounter{utterance}  
    & & & \multicolumn{4}{p{0.6\linewidth}}{
        \cellcolor[rgb]{0.9,0.9,0.9}{
            \makecell[{{p{\linewidth}}}]{
                \texttt{\tiny{[GM$\langle$P2]}}
                \texttt{حرّك: C, (2,4)} \\
            }
        }
    }
     \\ \\

    \theutterance \stepcounter{utterance}  
    & & & \multicolumn{2}{p{0.3\linewidth}}{
        \cellcolor[rgb]{0.9,0.9,0.9}{
            \makecell[{{p{\linewidth}}}]{
                \texttt{\tiny{[GM$|$GM]}}
                \texttt{حرّك: C, (2,4)} \\
            }
        }
    }
    & & \\ \\

    \theutterance \stepcounter{utterance}  
    & & & \multicolumn{2}{p{0.3\linewidth}}{
        \cellcolor[rgb]{0.9,0.9,0.9}{
            \makecell[{{p{\linewidth}}}]{
                \texttt{\tiny{[GM$|$GM]}}
                \texttt{عقوبة: (2,4) ليست فارغة، بل تحتوي على '─'.} \\
\\ 
\texttt{لقد جمعتما معاً 6 من أصل 8 عقوبة.} \\
\texttt{يرجى المحاولة مرة أخرى!} \\
            }
        }
    }
    & & \\ \\

    \theutterance \stepcounter{utterance}  
    & & \multicolumn{4}{p{0.6\linewidth}}{
        \cellcolor[rgb]{0.9,0.9,0.9}{
            \makecell[{{p{\linewidth}}}]{
                \texttt{\tiny{[GM$\rangle$P2]}}
                \texttt{عقوبة: (2,4) ليست فارغة، بل تحتوي على '─'.} \\
\\ 
\texttt{لقد جمعتما معاً 6 من أصل 8 عقوبة.} \\
\texttt{يرجى المحاولة مرة أخرى!} \\
            }
        }
    }
    & \\ \\

    \theutterance \stepcounter{utterance}  
    & & & \multicolumn{4}{p{0.6\linewidth}}{
        \cellcolor[rgb]{0.9,0.9,0.9}{
            \makecell[{{p{\linewidth}}}]{
                \texttt{\tiny{[GM$\langle$P2]}}
                \texttt{قل: أواجه مشكلة في تحديد الخلايا الفارغة بدقة. دعني أرسم لوحتي الحالية بوضوح:} \\
\\ 
\texttt{```} \\
\texttt{ 1234567} \\
\texttt{╔══╤══╤═╗} \\
\texttt{║◌P│◌◌│◌║ 1} \\
\texttt{╟──┼──┴─╢ 2} \\
\texttt{║◌◌│◌L◌◌║ 3} \\
\texttt{╟──┘◌◌◌◌║ 4} \\
\texttt{║◌◌◌◌◌◌C║ 5} \\
\texttt{║◌◌┌─┬──╢ 6} \\
\texttt{║◌◌│◌│◌◌║ 7} \\
\texttt{╚══╧═╧══╝} \\
\texttt{```} \\
\\ 
\texttt{الخلايا الفارغة (◌) هي:} \\
\texttt{(1,1), (3,1), (4,1), (5,1), (7,1)} \\
\texttt{(1,3), (2,3), (5,3), (6,3), (7,3)} \\
\texttt{(2,4), (3,4), (4,4), (5,4), (6,4), (7,4)} \\
\texttt{(1,5), (2,5), (3,5), (4,5), (5,5), (6,5)} \\
\texttt{(1,6), (2,6), (5,6)} \\
\texttt{(1,7), (2,7), (4,7), (5,7), (6,7), (7,7)} \\
\\ 
\texttt{هل يمكنك اختيار ثلاثة من هذه الخلايا تكون فارغة أيضاً على لوحتك؟} \\
            }
        }
    }
     \\ \\

    \theutterance \stepcounter{utterance}  
    & & & \multicolumn{2}{p{0.3\linewidth}}{
        \cellcolor[rgb]{0.9,0.9,0.9}{
            \makecell[{{p{\linewidth}}}]{
                \texttt{\tiny{[GM$|$GM]}}
                \texttt{قل: أواجه مشكلة في تحديد الخلايا الفارغة بدقة. دعني أرسم لوحتي الحالية بوضوح:} \\
\\ 
\texttt{```} \\
\texttt{ 1234567} \\
\texttt{╔══╤══╤═╗} \\
\texttt{║◌P│◌◌│◌║ 1} \\
\texttt{╟──┼──┴─╢ 2} \\
\texttt{║◌◌│◌L◌◌║ 3} \\
\texttt{╟──┘◌◌◌◌║ 4} \\
\texttt{║◌◌◌◌◌◌C║ 5} \\
\texttt{║◌◌┌─┬──╢ 6} \\
\texttt{║◌◌│◌│◌◌║ 7} \\
\texttt{╚══╧═╧══╝} \\
\texttt{```} \\
\\ 
\texttt{الخلايا الفارغة (◌) هي:} \\
\texttt{(1,1), (3,1), (4,1), (5,1), (7,1)} \\
\texttt{(1,3), (2,3), (5,3), (6,3), (7,3)} \\
\texttt{(2,4), (3,4), (4,4), (5,4), (6,4), (7,4)} \\
\texttt{(1,5), (2,5), (3,5), (4,5), (5,5), (6,5)} \\
\texttt{(1,6), (2,6), (5,6)} \\
\texttt{(1,7), (2,7), (4,7), (5,7), (6,7), (7,7)} \\
\\ 
\texttt{هل يمكنك اختيار ثلاثة من هذه الخلايا تكون فارغة أيضاً على لوحتك؟} \\
            }
        }
    }
    & & \\ \\

    \theutterance \stepcounter{utterance}  
    & & & \multicolumn{2}{p{0.3\linewidth}}{
        \cellcolor[rgb]{0.9,0.9,0.9}{
            \makecell[{{p{\linewidth}}}]{
                \texttt{\tiny{[GM$|$GM]}}
                \texttt{Invalid format: tail is not empty: '} \\
\\ 
\\ 
\texttt{ 1234567} \\
\texttt{╔══╤══╤═╗} \\
\texttt{║◌P│◌◌│◌║ 1} \\
\texttt{╟──┼──┴─╢ 2} \\
\texttt{║◌◌│◌L◌◌║ 3} \\
\texttt{╟──┘◌◌◌◌║ 4} \\
\texttt{║◌◌◌◌◌◌C║ 5} \\
\texttt{║◌◌┌─┬──╢ 6} \\
\texttt{║◌◌│◌│◌◌║ 7} \\
\texttt{╚══╧═╧══╝} \\
\\ 
\\ 
\texttt{الخلايا الفارغة (◌) هي:} \\
\texttt{(1,1), (3,1), (4,1), (5,1), (7,1)} \\
\texttt{(1,3), (2,3), (5,3), (6,3), (7,3)} \\
\texttt{(2,4), (3,4), (4,4), (5,4), (6,4), (7,4)} \\
\texttt{(1,5), (2,5), (3,5), (4,5), (5,5), (6,5)} \\
\texttt{(1,6), (2,6), (5,6)} \\
\texttt{(1,7), (2,7), (4,7), (5,7), (6,7), (7,7)} \\
\\ 
\texttt{هل يمكنك اختيار ثلاثة من هذه الخلايا تكون فارغة أيضاً على لوحتك؟} \\
            }
        }
    }
    & & \\ \\

    \theutterance \stepcounter{utterance}  
    & & \multicolumn{4}{p{0.6\linewidth}}{
        \cellcolor[rgb]{0.9,0.9,0.9}{
            \makecell[{{p{\linewidth}}}]{
                \texttt{\tiny{[GM$\rangle$P2]}}
                \texttt{عقوبة: يجب ألا تحتوي رسالتك على أي شيء بعد الأمر!} \\
\\ 
\texttt{تأكد من أن ردك يحتوي فقط على `قل: <رسالة>` أو `حرّك: <كائن>, (<X>, <Y>)`، ولا شيء غير ذلك!} \\
\\ 
\texttt{لقد جمعتما معاً 7 من أصل 8 عقوبة.} \\
\texttt{يرجى المحاولة مرة أخرى!} \\
            }
        }
    }
    & \\ \\

    \theutterance \stepcounter{utterance}  
    & & & \multicolumn{4}{p{0.6\linewidth}}{
        \cellcolor[rgb]{0.9,0.9,0.9}{
            \makecell[{{p{\linewidth}}}]{
                \texttt{\tiny{[GM$\langle$P2]}}
                \texttt{قل: أواجه مشكلة في تحديد الخلايا الفارغة بدقة. الخلايا الفارغة المؤكدة لدي هي: (1,1) و(3,1) و(4,1) و(5,1) و(7,1) و(1,3) و(2,3) و(5,3) و(6,3) و(7,3) و(2,4) و(3,4) و(4,4) و(5,4) و(6,4) و(7,4) و(1,5) و(2,5) و(3,5) و(4,5) و(5,5) و(6,5) و(1,6) و(2,6) و(5,6) و(1,7) و(2,7) و(4,7) و(5,7) و(6,7) و(7,7). هل يمكنك اختيار ثلاثة منها فارغة لديك أيضاً؟} \\
            }
        }
    }
     \\ \\

    \theutterance \stepcounter{utterance}  
    & & & \multicolumn{2}{p{0.3\linewidth}}{
        \cellcolor[rgb]{0.9,0.9,0.9}{
            \makecell[{{p{\linewidth}}}]{
                \texttt{\tiny{[GM$|$GM]}}
                \texttt{قل: أواجه مشكلة في تحديد الخلايا الفارغة بدقة. الخلايا الفارغة المؤكدة لدي هي: (1,1) و(3,1) و(4,1) و(5,1) و(7,1) و(1,3) و(2,3) و(5,3) و(6,3) و(7,3) و(2,4) و(3,4) و(4,4) و(5,4) و(6,4) و(7,4) و(1,5) و(2,5) و(3,5) و(4,5) و(5,5) و(6,5) و(1,6) و(2,6) و(5,6) و(1,7) و(2,7) و(4,7) و(5,7) و(6,7) و(7,7). هل يمكنك اختيار ثلاثة منها فارغة لديك أيضاً؟} \\
            }
        }
    }
    & & \\ \\

    \theutterance \stepcounter{utterance}  
    & & \multicolumn{4}{p{0.6\linewidth}}{
        \cellcolor[rgb]{0.9,0.9,0.9}{
            \makecell[{{p{\linewidth}}}]{
                \texttt{\tiny{[P1$\langle$GM]}}
                \texttt{تم تمرير رسالتك إلى اللاعب الآخر.} \\
\\ 
\\ 
\texttt{أنت تلعب حالياً الجولة 4 من أصل 12 كحد أقصى.} \\
\texttt{لقد جمعتما معاً 7 من أصل 8 عقوبة.} \\
\texttt{أرسل اللاعب الآخر هذه الرسالة:} \\
\\ 
\\ 
\texttt{"أواجه مشكلة في تحديد الخلايا الفارغة بدقة. الخلايا الفارغة المؤكدة لدي هي: (1,1) و(3,1) و(4,1) و(5,1) و(7,1) و(1,3) و(2,3) و(5,3) و(6,3) و(7,3) و(2,4) و(3,4) و(4,4) و(5,4) و(6,4) و(7,4) و(1,5) و(2,5) و(3,5) و(4,5) و(5,5) و(6,5) و(1,6) و(2,6) و(5,6) و(1,7) و(2,7) و(4,7) و(5,7) و(6,7) و(7,7). هل يمكنك اختيار ثلاثة منها فارغة لديك أيضاً؟"} \\
\\ 
\\ 
\texttt{ما هو أمرك التالي؟} \\
            }
        }
    }
    & \\ \\

    \theutterance \stepcounter{utterance}  
    & \multicolumn{4}{p{0.6\linewidth}}{
        \cellcolor[rgb]{0.9,0.9,0.9}{
            \makecell[{{p{\linewidth}}}]{
                \texttt{\tiny{[P1$\rangle$GM]}}
                \texttt{قل: لنتحقق من الاقتراحات السابقة أولاً. (3,2) و (5,4) و (6,5) كلها فارغة على لوحتي. هل يمكنك التأكد مما إذا كانت هذه الخلايا الثلاث فارغة على لوحتك أيضاً؟ إذا كانت كذلك، فلنستخدمها. وإلا سأقترح بدائل من قائمتك.} \\
            }
        }
    }
    & & \\ \\

    \theutterance \stepcounter{utterance}  
    & & & \multicolumn{2}{p{0.3\linewidth}}{
        \cellcolor[rgb]{0.9,0.9,0.9}{
            \makecell[{{p{\linewidth}}}]{
                \texttt{\tiny{[GM$|$GM]}}
                \texttt{قل: لنتحقق من الاقتراحات السابقة أولاً. (3,2) و (5,4) و (6,5) كلها فارغة على لوحتي. هل يمكنك التأكد مما إذا كانت هذه الخلايا الثلاث فارغة على لوحتك أيضاً؟ إذا كانت كذلك، فلنستخدمها. وإلا سأقترح بدائل من قائمتك.} \\
            }
        }
    }
    & & \\ \\

    \theutterance \stepcounter{utterance}  
    & & \multicolumn{4}{p{0.6\linewidth}}{
        \cellcolor[rgb]{0.9,0.9,0.9}{
            \makecell[{{p{\linewidth}}}]{
                \texttt{\tiny{[GM$\rangle$P2]}}
                \texttt{تم تمرير رسالتك إلى اللاعب الآخر.} \\
\\ 
\\ 
\texttt{أنت تلعب حالياً الجولة 5 من أصل 12 كحد أقصى.} \\
\texttt{لقد جمعتما معاً 7 من أصل 8 عقوبة.} \\
\texttt{أرسل اللاعب الآخر هذه الرسالة:} \\
\\ 
\\ 
\texttt{"لنتحقق من الاقتراحات السابقة أولاً. (3,2) و (5,4) و (6,5) كلها فارغة على لوحتي. هل يمكنك التأكد مما إذا كانت هذه الخلايا الثلاث فارغة على لوحتك أيضاً؟ إذا كانت كذلك، فلنستخدمها. وإلا سأقترح بدائل من قائمتك."} \\
\\ 
\\ 
\texttt{ما هو أمرك التالي؟} \\
            }
        }
    }
    & \\ \\

    \theutterance \stepcounter{utterance}  
    & & & \multicolumn{4}{p{0.6\linewidth}}{
        \cellcolor[rgb]{0.9,0.9,0.9}{
            \makecell[{{p{\linewidth}}}]{
                \texttt{\tiny{[GM$\langle$P2]}}
                \texttt{قل: (3,2) ليست فارغة على لوحتي، بل تحتوي على '┼'. وكذلك (2,4) تحتوي على '─' لذا لم تنجح محاولاتي السابقة. لكن (5,4) و(6,5) فارغتان. هل تريد استخدام (5,4) و(6,5) وإيجاد بديل للثالثة؟} \\
            }
        }
    }
     \\ \\

    \theutterance \stepcounter{utterance}  
    & & & \multicolumn{2}{p{0.3\linewidth}}{
        \cellcolor[rgb]{0.9,0.9,0.9}{
            \makecell[{{p{\linewidth}}}]{
                \texttt{\tiny{[GM$|$GM]}}
                \texttt{قل: (3,2) ليست فارغة على لوحتي، بل تحتوي على '┼'. وكذلك (2,4) تحتوي على '─' لذا لم تنجح محاولاتي السابقة. لكن (5,4) و(6,5) فارغتان. هل تريد استخدام (5,4) و(6,5) وإيجاد بديل للثالثة؟} \\
            }
        }
    }
    & & \\ \\

    \theutterance \stepcounter{utterance}  
    & & \multicolumn{4}{p{0.6\linewidth}}{
        \cellcolor[rgb]{0.9,0.9,0.9}{
            \makecell[{{p{\linewidth}}}]{
                \texttt{\tiny{[P1$\langle$GM]}}
                \texttt{تم تمرير رسالتك إلى اللاعب الآخر.} \\
\\ 
\\ 
\texttt{أنت تلعب حالياً الجولة 5 من أصل 12 كحد أقصى.} \\
\texttt{لقد جمعتما معاً 7 من أصل 8 عقوبة.} \\
\texttt{أرسل اللاعب الآخر هذه الرسالة:} \\
\\ 
\\ 
\texttt{"(3,2) ليست فارغة على لوحتي، بل تحتوي على '┼'. وكذلك (2,4) تحتوي على '─' لذا لم تنجح محاولاتي السابقة. لكن (5,4) و(6,5) فارغتان. هل تريد استخدام (5,4) و(6,5) وإيجاد بديل للثالثة؟"} \\
\\ 
\\ 
\texttt{ما هو أمرك التالي؟} \\
            }
        }
    }
    & \\ \\

    \theutterance \stepcounter{utterance}  
    & \multicolumn{4}{p{0.6\linewidth}}{
        \cellcolor[rgb]{0.9,0.9,0.9}{
            \makecell[{{p{\linewidth}}}]{
                \texttt{\tiny{[P1$\rangle$GM]}}
                \texttt{قل: حسناً، لنستخدم (5,4) و (6,5). بالنسبة للخيار الثالث، من قائمتك، لنتحقق من (3,6) و (4,6) و (5,6). هل أي منها فارغ على لوحتك؟ على لوحتي جميعها فارغة.} \\
            }
        }
    }
    & & \\ \\

    \theutterance \stepcounter{utterance}  
    & & & \multicolumn{2}{p{0.3\linewidth}}{
        \cellcolor[rgb]{0.9,0.9,0.9}{
            \makecell[{{p{\linewidth}}}]{
                \texttt{\tiny{[GM$|$GM]}}
                \texttt{قل: حسناً، لنستخدم (5,4) و (6,5). بالنسبة للخيار الثالث، من قائمتك، لنتحقق من (3,6) و (4,6) و (5,6). هل أي منها فارغ على لوحتك؟ على لوحتي جميعها فارغة.} \\
            }
        }
    }
    & & \\ \\

    \theutterance \stepcounter{utterance}  
    & & \multicolumn{4}{p{0.6\linewidth}}{
        \cellcolor[rgb]{0.9,0.9,0.9}{
            \makecell[{{p{\linewidth}}}]{
                \texttt{\tiny{[GM$\rangle$P2]}}
                \texttt{تم تمرير رسالتك إلى اللاعب الآخر.} \\
\\ 
\\ 
\texttt{أنت تلعب حالياً الجولة 6 من أصل 12 كحد أقصى.} \\
\texttt{لقد جمعتما معاً 7 من أصل 8 عقوبة.} \\
\texttt{أرسل اللاعب الآخر هذه الرسالة:} \\
\\ 
\\ 
\texttt{"حسناً، لنستخدم (5,4) و (6,5). بالنسبة للخيار الثالث، من قائمتك، لنتحقق من (3,6) و (4,6) و (5,6). هل أي منها فارغ على لوحتك؟ على لوحتي جميعها فارغة."} \\
\\ 
\\ 
\texttt{ما هو أمرك التالي؟} \\
            }
        }
    }
    & \\ \\

    \theutterance \stepcounter{utterance}  
    & & & \multicolumn{4}{p{0.6\linewidth}}{
        \cellcolor[rgb]{0.9,0.9,0.9}{
            \makecell[{{p{\linewidth}}}]{
                \texttt{\tiny{[GM$\langle$P2]}}
                \texttt{قل: (3,6) و(4,6) و(5,6) جميعها تحتوي على إطارات أو خطوط على لوحتي وليست فارغة. الخلايا الفارغة من الصف 6 لدي هي فقط (1,6) و(2,6) و(5,6) {-} انتظر، (5,6) ليست فارغة بل تحتوي على '┬'. إذاً (1,6) و(2,6) فقط. هل هما فارغتان على لوحتك؟} \\
            }
        }
    }
     \\ \\

    \theutterance \stepcounter{utterance}  
    & & & \multicolumn{2}{p{0.3\linewidth}}{
        \cellcolor[rgb]{0.9,0.9,0.9}{
            \makecell[{{p{\linewidth}}}]{
                \texttt{\tiny{[GM$|$GM]}}
                \texttt{قل: (3,6) و(4,6) و(5,6) جميعها تحتوي على إطارات أو خطوط على لوحتي وليست فارغة. الخلايا الفارغة من الصف 6 لدي هي فقط (1,6) و(2,6) و(5,6) {-} انتظر، (5,6) ليست فارغة بل تحتوي على '┬'. إذاً (1,6) و(2,6) فقط. هل هما فارغتان على لوحتك؟} \\
            }
        }
    }
    & & \\ \\

    \theutterance \stepcounter{utterance}  
    & & \multicolumn{4}{p{0.6\linewidth}}{
        \cellcolor[rgb]{0.9,0.9,0.9}{
            \makecell[{{p{\linewidth}}}]{
                \texttt{\tiny{[P1$\langle$GM]}}
                \texttt{تم تمرير رسالتك إلى اللاعب الآخر.} \\
\\ 
\\ 
\texttt{أنت تلعب حالياً الجولة 6 من أصل 12 كحد أقصى.} \\
\texttt{لقد جمعتما معاً 7 من أصل 8 عقوبة.} \\
\texttt{أرسل اللاعب الآخر هذه الرسالة:} \\
\\ 
\\ 
\texttt{"(3,6) و(4,6) و(5,6) جميعها تحتوي على إطارات أو خطوط على لوحتي وليست فارغة. الخلايا الفارغة من الصف 6 لدي هي فقط (1,6) و(2,6) و(5,6) {-} انتظر، (5,6) ليست فارغة بل تحتوي على '┬'. إذاً (1,6) و(2,6) فقط. هل هما فارغتان على لوحتك؟"} \\
\\ 
\\ 
\texttt{ما هو أمرك التالي؟} \\
            }
        }
    }
    & \\ \\

    \theutterance \stepcounter{utterance}  
    & \multicolumn{4}{p{0.6\linewidth}}{
        \cellcolor[rgb]{0.9,0.9,0.9}{
            \makecell[{{p{\linewidth}}}]{
                \texttt{\tiny{[P1$\rangle$GM]}}
                \texttt{حرّك: C, (5,4)} \\
            }
        }
    }
    & & \\ \\

    \theutterance \stepcounter{utterance}  
    & & & \multicolumn{2}{p{0.3\linewidth}}{
        \cellcolor[rgb]{0.9,0.9,0.9}{
            \makecell[{{p{\linewidth}}}]{
                \texttt{\tiny{[GM$|$GM]}}
                \texttt{حرّك: C, (5,4)} \\
            }
        }
    }
    & & \\ \\

    \theutterance \stepcounter{utterance}  
    & & & \multicolumn{2}{p{0.3\linewidth}}{
        \cellcolor[rgb]{0.9,0.9,0.9}{
            \makecell[{{p{\linewidth}}}]{
                \texttt{\tiny{[GM$|$GM]}}
                \texttt{C moved to (5, 4)} \\
            }
        }
    }
    & & \\ \\

    \theutterance \stepcounter{utterance}  
    & & \multicolumn{4}{p{0.6\linewidth}}{
        \cellcolor[rgb]{0.9,0.9,0.9}{
            \makecell[{{p{\linewidth}}}]{
                \texttt{\tiny{[GM$\rangle$P2]}}
                \texttt{تم تمرير رسالتك إلى اللاعب الآخر.} \\
\\ 
\\ 
\texttt{أنت تلعب حالياً الجولة 7 من أصل 12 كحد أقصى.} \\
\texttt{لقد جمعتما معاً 7 من أصل 8 عقوبة.} \\
\texttt{حرّك اللاعب الآخر كائناً على شبكته.} \\
\\ 
\\ 
\texttt{ما هو أمرك التالي؟} \\
            }
        }
    }
    & \\ \\

    \theutterance \stepcounter{utterance}  
    & & & \multicolumn{4}{p{0.6\linewidth}}{
        \cellcolor[rgb]{0.9,0.9,0.9}{
            \makecell[{{p{\linewidth}}}]{
                \texttt{\tiny{[GM$\langle$P2]}}
                \texttt{حرّك: P, (6,5)} \\
            }
        }
    }
     \\ \\

    \theutterance \stepcounter{utterance}  
    & & & \multicolumn{2}{p{0.3\linewidth}}{
        \cellcolor[rgb]{0.9,0.9,0.9}{
            \makecell[{{p{\linewidth}}}]{
                \texttt{\tiny{[GM$|$GM]}}
                \texttt{حرّك: P, (6,5)} \\
            }
        }
    }
    & & \\ \\

    \theutterance \stepcounter{utterance}  
    & & & \multicolumn{2}{p{0.3\linewidth}}{
        \cellcolor[rgb]{0.9,0.9,0.9}{
            \makecell[{{p{\linewidth}}}]{
                \texttt{\tiny{[GM$|$GM]}}
                \texttt{P moved to (6, 5)} \\
            }
        }
    }
    & & \\ \\

    \theutterance \stepcounter{utterance}  
    & & \multicolumn{4}{p{0.6\linewidth}}{
        \cellcolor[rgb]{0.9,0.9,0.9}{
            \makecell[{{p{\linewidth}}}]{
                \texttt{\tiny{[P1$\langle$GM]}}
                \texttt{تم تحريك 'C' إلى (5,4) بنجاح. تبدو شبكتك المحدّثة هكذا:} \\
\texttt{```} \\
\texttt{ 1234567} \\
\texttt{╔══╤══╤═╗} \\
\texttt{║◌◌│◌◌│◌║ 1} \\
\texttt{╟──┼──┴─╢ 2} \\
\texttt{║L◌│◌◌◌◌║ 3} \\
\texttt{╟──┘◌C◌P║ 4} \\
\texttt{║◌◌◌◌◌◌◌║ 5} \\
\texttt{║◌◌┌─┬──╢ 6} \\
\texttt{║◌◌│◌│◌◌║ 7} \\
\texttt{╚══╧═╧══╝} \\
\texttt{```} \\
\texttt{أنت تلعب حالياً الجولة 7 من أصل 12 كحد أقصى.} \\
\texttt{لقد جمعتما معاً 7 من أصل 8 عقوبة.} \\
\texttt{حرّك اللاعب الآخر كائناً على شبكته.} \\
\\ 
\\ 
\texttt{ما هو أمرك التالي؟} \\
            }
        }
    }
    & \\ \\

    \theutterance \stepcounter{utterance}  
    & \multicolumn{4}{p{0.6\linewidth}}{
        \cellcolor[rgb]{0.9,0.9,0.9}{
            \makecell[{{p{\linewidth}}}]{
                \texttt{\tiny{[P1$\rangle$GM]}}
                \texttt{قل: قمت بتحريك C إلى (5,4). (1,6) و (2,6) فارغتان على لوحتي. فلنستخدم إحداهما. أقترح (2,6) للكائن L. هل (2,6) فارغة على لوحتك؟} \\
            }
        }
    }
    & & \\ \\

    \theutterance \stepcounter{utterance}  
    & & & \multicolumn{2}{p{0.3\linewidth}}{
        \cellcolor[rgb]{0.9,0.9,0.9}{
            \makecell[{{p{\linewidth}}}]{
                \texttt{\tiny{[GM$|$GM]}}
                \texttt{قل: قمت بتحريك C إلى (5,4). (1,6) و (2,6) فارغتان على لوحتي. فلنستخدم إحداهما. أقترح (2,6) للكائن L. هل (2,6) فارغة على لوحتك؟} \\
            }
        }
    }
    & & \\ \\

    \theutterance \stepcounter{utterance}  
    & & \multicolumn{4}{p{0.6\linewidth}}{
        \cellcolor[rgb]{0.9,0.9,0.9}{
            \makecell[{{p{\linewidth}}}]{
                \texttt{\tiny{[GM$\rangle$P2]}}
                \texttt{تم تحريك 'P' إلى (6,5) بنجاح. تبدو شبكتك المحدّثة هكذا:} \\
\texttt{```} \\
\texttt{ 1234567} \\
\texttt{╔══╤══╤═╗} \\
\texttt{║◌◌│◌◌│◌║ 1} \\
\texttt{╟──┼──┴─╢ 2} \\
\texttt{║◌◌│◌L◌◌║ 3} \\
\texttt{╟──┘◌◌◌◌║ 4} \\
\texttt{║◌◌◌◌◌PC║ 5} \\
\texttt{║◌◌┌─┬──╢ 6} \\
\texttt{║◌◌│◌│◌◌║ 7} \\
\texttt{╚══╧═╧══╝} \\
\texttt{```} \\
\texttt{أنت تلعب حالياً الجولة 8 من أصل 12 كحد أقصى.} \\
\texttt{لقد جمعتما معاً 7 من أصل 8 عقوبة.} \\
\texttt{أرسل اللاعب الآخر هذه الرسالة:} \\
\\ 
\\ 
\texttt{"قمت بتحريك C إلى (5,4). (1,6) و (2,6) فارغتان على لوحتي. فلنستخدم إحداهما. أقترح (2,6) للكائن L. هل (2,6) فارغة على لوحتك؟"} \\
\\ 
\\ 
\texttt{ما هو أمرك التالي؟} \\
            }
        }
    }
    & \\ \\

    \theutterance \stepcounter{utterance}  
    & & & \multicolumn{4}{p{0.6\linewidth}}{
        \cellcolor[rgb]{0.9,0.9,0.9}{
            \makecell[{{p{\linewidth}}}]{
                \texttt{\tiny{[GM$\langle$P2]}}
                \texttt{حرّك: L, (2,6)} \\
            }
        }
    }
     \\ \\

    \theutterance \stepcounter{utterance}  
    & & & \multicolumn{2}{p{0.3\linewidth}}{
        \cellcolor[rgb]{0.9,0.9,0.9}{
            \makecell[{{p{\linewidth}}}]{
                \texttt{\tiny{[GM$|$GM]}}
                \texttt{حرّك: L, (2,6)} \\
            }
        }
    }
    & & \\ \\

    \theutterance \stepcounter{utterance}  
    & & & \multicolumn{2}{p{0.3\linewidth}}{
        \cellcolor[rgb]{0.9,0.9,0.9}{
            \makecell[{{p{\linewidth}}}]{
                \texttt{\tiny{[GM$|$GM]}}
                \texttt{L moved to (2, 6)} \\
            }
        }
    }
    & & \\ \\

    \theutterance \stepcounter{utterance}  
    & & \multicolumn{4}{p{0.6\linewidth}}{
        \cellcolor[rgb]{0.9,0.9,0.9}{
            \makecell[{{p{\linewidth}}}]{
                \texttt{\tiny{[P1$\langle$GM]}}
                \texttt{تم تمرير رسالتك إلى اللاعب الآخر.} \\
\\ 
\\ 
\texttt{أنت تلعب حالياً الجولة 8 من أصل 12 كحد أقصى.} \\
\texttt{لقد جمعتما معاً 7 من أصل 8 عقوبة.} \\
\texttt{حرّك اللاعب الآخر كائناً على شبكته.} \\
\\ 
\\ 
\texttt{ما هو أمرك التالي؟} \\
            }
        }
    }
    & \\ \\

    \theutterance \stepcounter{utterance}  
    & \multicolumn{4}{p{0.6\linewidth}}{
        \cellcolor[rgb]{0.9,0.9,0.9}{
            \makecell[{{p{\linewidth}}}]{
                \texttt{\tiny{[P1$\rangle$GM]}}
                \texttt{حرّك: L, (2,6)} \\
            }
        }
    }
    & & \\ \\

    \theutterance \stepcounter{utterance}  
    & & & \multicolumn{2}{p{0.3\linewidth}}{
        \cellcolor[rgb]{0.9,0.9,0.9}{
            \makecell[{{p{\linewidth}}}]{
                \texttt{\tiny{[GM$|$GM]}}
                \texttt{حرّك: L, (2,6)} \\
            }
        }
    }
    & & \\ \\

    \theutterance \stepcounter{utterance}  
    & & & \multicolumn{2}{p{0.3\linewidth}}{
        \cellcolor[rgb]{0.9,0.9,0.9}{
            \makecell[{{p{\linewidth}}}]{
                \texttt{\tiny{[GM$|$GM]}}
                \texttt{L moved to (2, 6)} \\
            }
        }
    }
    & & \\ \\

    \theutterance \stepcounter{utterance}  
    & & \multicolumn{4}{p{0.6\linewidth}}{
        \cellcolor[rgb]{0.9,0.9,0.9}{
            \makecell[{{p{\linewidth}}}]{
                \texttt{\tiny{[GM$\rangle$P2]}}
                \texttt{تم تحريك 'L' إلى (2,6) بنجاح. تبدو شبكتك المحدّثة هكذا:} \\
\texttt{```} \\
\texttt{ 1234567} \\
\texttt{╔══╤══╤═╗} \\
\texttt{║◌◌│◌◌│◌║ 1} \\
\texttt{╟──┼──┴─╢ 2} \\
\texttt{║◌◌│◌◌◌◌║ 3} \\
\texttt{╟──┘◌◌◌◌║ 4} \\
\texttt{║◌◌◌◌◌PC║ 5} \\
\texttt{║◌L┌─┬──╢ 6} \\
\texttt{║◌◌│◌│◌◌║ 7} \\
\texttt{╚══╧═╧══╝} \\
\texttt{```} \\
\texttt{أنت تلعب حالياً الجولة 9 من أصل 12 كحد أقصى.} \\
\texttt{لقد جمعتما معاً 7 من أصل 8 عقوبة.} \\
\texttt{حرّك اللاعب الآخر كائناً على شبكته.} \\
\\ 
\\ 
\texttt{ما هو أمرك التالي؟} \\
            }
        }
    }
    & \\ \\

    \theutterance \stepcounter{utterance}  
    & & & \multicolumn{4}{p{0.6\linewidth}}{
        \cellcolor[rgb]{0.9,0.9,0.9}{
            \makecell[{{p{\linewidth}}}]{
                \texttt{\tiny{[GM$\langle$P2]}}
                \texttt{حرّك: C, (5,4)} \\
            }
        }
    }
     \\ \\

    \theutterance \stepcounter{utterance}  
    & & & \multicolumn{2}{p{0.3\linewidth}}{
        \cellcolor[rgb]{0.9,0.9,0.9}{
            \makecell[{{p{\linewidth}}}]{
                \texttt{\tiny{[GM$|$GM]}}
                \texttt{حرّك: C, (5,4)} \\
            }
        }
    }
    & & \\ \\

    \theutterance \stepcounter{utterance}  
    & & & \multicolumn{2}{p{0.3\linewidth}}{
        \cellcolor[rgb]{0.9,0.9,0.9}{
            \makecell[{{p{\linewidth}}}]{
                \texttt{\tiny{[GM$|$GM]}}
                \texttt{C moved to (5, 4)} \\
            }
        }
    }
    & & \\ \\

    \theutterance \stepcounter{utterance}  
    & & \multicolumn{4}{p{0.6\linewidth}}{
        \cellcolor[rgb]{0.9,0.9,0.9}{
            \makecell[{{p{\linewidth}}}]{
                \texttt{\tiny{[P1$\langle$GM]}}
                \texttt{تم تحريك 'L' إلى (2,6) بنجاح. تبدو شبكتك المحدّثة هكذا:} \\
\texttt{```} \\
\texttt{ 1234567} \\
\texttt{╔══╤══╤═╗} \\
\texttt{║◌◌│◌◌│◌║ 1} \\
\texttt{╟──┼──┴─╢ 2} \\
\texttt{║◌◌│◌◌◌◌║ 3} \\
\texttt{╟──┘◌C◌P║ 4} \\
\texttt{║◌◌◌◌◌◌◌║ 5} \\
\texttt{║◌L┌─┬──╢ 6} \\
\texttt{║◌◌│◌│◌◌║ 7} \\
\texttt{╚══╧═╧══╝} \\
\texttt{```} \\
\texttt{أنت تلعب حالياً الجولة 9 من أصل 12 كحد أقصى.} \\
\texttt{لقد جمعتما معاً 7 من أصل 8 عقوبة.} \\
\texttt{حرّك اللاعب الآخر كائناً على شبكته.} \\
\\ 
\\ 
\texttt{ما هو أمرك التالي؟} \\
            }
        }
    }
    & \\ \\

    \theutterance \stepcounter{utterance}  
    & \multicolumn{4}{p{0.6\linewidth}}{
        \cellcolor[rgb]{0.9,0.9,0.9}{
            \makecell[{{p{\linewidth}}}]{
                \texttt{\tiny{[P1$\rangle$GM]}}
                \texttt{قل: قمت بتحريك L إلى (2,6). الآن بقي P للتحريك إلى (6,5). هل (6,5) ما زالت فارغة على لوحتك؟} \\
            }
        }
    }
    & & \\ \\

    \theutterance \stepcounter{utterance}  
    & & & \multicolumn{2}{p{0.3\linewidth}}{
        \cellcolor[rgb]{0.9,0.9,0.9}{
            \makecell[{{p{\linewidth}}}]{
                \texttt{\tiny{[GM$|$GM]}}
                \texttt{قل: قمت بتحريك L إلى (2,6). الآن بقي P للتحريك إلى (6,5). هل (6,5) ما زالت فارغة على لوحتك؟} \\
            }
        }
    }
    & & \\ \\

    \theutterance \stepcounter{utterance}  
    & & \multicolumn{4}{p{0.6\linewidth}}{
        \cellcolor[rgb]{0.9,0.9,0.9}{
            \makecell[{{p{\linewidth}}}]{
                \texttt{\tiny{[GM$\rangle$P2]}}
                \texttt{تم تحريك 'C' إلى (5,4) بنجاح. تبدو شبكتك المحدّثة هكذا:} \\
\texttt{```} \\
\texttt{ 1234567} \\
\texttt{╔══╤══╤═╗} \\
\texttt{║◌◌│◌◌│◌║ 1} \\
\texttt{╟──┼──┴─╢ 2} \\
\texttt{║◌◌│◌◌◌◌║ 3} \\
\texttt{╟──┘◌C◌◌║ 4} \\
\texttt{║◌◌◌◌◌P◌║ 5} \\
\texttt{║◌L┌─┬──╢ 6} \\
\texttt{║◌◌│◌│◌◌║ 7} \\
\texttt{╚══╧═╧══╝} \\
\texttt{```} \\
\texttt{أنت تلعب حالياً الجولة 10 من أصل 12 كحد أقصى.} \\
\texttt{لقد جمعتما معاً 7 من أصل 8 عقوبة.} \\
\texttt{أرسل اللاعب الآخر هذه الرسالة:} \\
\\ 
\\ 
\texttt{"قمت بتحريك L إلى (2,6). الآن بقي P للتحريك إلى (6,5). هل (6,5) ما زالت فارغة على لوحتك؟"} \\
\\ 
\\ 
\texttt{ما هو أمرك التالي؟} \\
            }
        }
    }
    & \\ \\

    \theutterance \stepcounter{utterance}  
    & & & \multicolumn{4}{p{0.6\linewidth}}{
        \cellcolor[rgb]{0.9,0.9,0.9}{
            \makecell[{{p{\linewidth}}}]{
                \texttt{\tiny{[GM$\langle$P2]}}
                \texttt{P موجود بالفعل في (6,5) على لوحتي! جميع كائناتي الآن في مواقعها: C عند (5,4)، L عند (2,6)، و P عند (6,5). هل كائناتك أيضاً في هذه المواقع؟ إذا كان الأمر كذلك، يمكننا إنهاء اللعبة.} \\
            }
        }
    }
     \\ \\

    \theutterance \stepcounter{utterance}  
    & & & \multicolumn{2}{p{0.3\linewidth}}{
        \cellcolor[rgb]{0.9,0.9,0.9}{
            \makecell[{{p{\linewidth}}}]{
                \texttt{\tiny{[GM$|$GM]}}
                \texttt{P موجود بالفعل في (6,5) على لوحتي! جميع كائناتي الآن في مواقعها: C عند (5,4)، L عند (2,6)، و P عند (6,5). هل كائناتك أيضاً في هذه المواقع؟ إذا كان الأمر كذلك، يمكننا إنهاء اللعبة.} \\
            }
        }
    }
    & & \\ \\

    \theutterance \stepcounter{utterance}  
    & & & \multicolumn{2}{p{0.3\linewidth}}{
        \cellcolor[rgb]{0.9,0.9,0.9}{
            \makecell[{{p{\linewidth}}}]{
                \texttt{\tiny{[GM$|$GM]}}
                \texttt{Invalid response format} \\
            }
        }
    }
    & & \\ \\

    \theutterance \stepcounter{utterance}  
    & & \multicolumn{4}{p{0.6\linewidth}}{
        \cellcolor[rgb]{0.9,0.9,0.9}{
            \makecell[{{p{\linewidth}}}]{
                \texttt{\tiny{[GM$\rangle$P2]}}
                \texttt{عقوبة: رسالتك ليست بالتنسيق المتوقع!} \\
\\ 
\texttt{تأكد من أن ردك يحتوي فقط على `قل: <رسالة>` أو `حرّك: <كائن>, (<X>, <Y>)`، ولا شيء غير ذلك!} \\
\\ 
\texttt{لقد جمعتما معاً 8 من أصل 8 عقوبة.} \\
\texttt{يرجى المحاولة مرة أخرى!} \\
            }
        }
    }
    & \\ \\

    \theutterance \stepcounter{utterance}  
    & & & \multicolumn{4}{p{0.6\linewidth}}{
        \cellcolor[rgb]{0.9,0.9,0.9}{
            \makecell[{{p{\linewidth}}}]{
                \texttt{\tiny{[GM$\langle$P2]}}
                \texttt{قل: P موجود بالفعل في (6,5) على لوحتي. جميع كائناتي الآن في مواقعها: C عند (5,4)، L عند (2,6)، و P عند (6,5). هل كائناتك أيضاً في هذه المواقع؟ إذا كان الأمر كذلك، يمكننا إنهاء اللعبة.} \\
            }
        }
    }
     \\ \\

    \theutterance \stepcounter{utterance}  
    & & & \multicolumn{2}{p{0.3\linewidth}}{
        \cellcolor[rgb]{0.9,0.9,0.9}{
            \makecell[{{p{\linewidth}}}]{
                \texttt{\tiny{[GM$|$GM]}}
                \texttt{قل: P موجود بالفعل في (6,5) على لوحتي. جميع كائناتي الآن في مواقعها: C عند (5,4)، L عند (2,6)، و P عند (6,5). هل كائناتك أيضاً في هذه المواقع؟ إذا كان الأمر كذلك، يمكننا إنهاء اللعبة.} \\
            }
        }
    }
    & & \\ \\

    \theutterance \stepcounter{utterance}  
    & & \multicolumn{4}{p{0.6\linewidth}}{
        \cellcolor[rgb]{0.9,0.9,0.9}{
            \makecell[{{p{\linewidth}}}]{
                \texttt{\tiny{[P1$\langle$GM]}}
                \texttt{تم تمرير رسالتك إلى اللاعب الآخر.} \\
\\ 
\\ 
\texttt{أنت تلعب حالياً الجولة 10 من أصل 12 كحد أقصى.} \\
\texttt{لقد جمعتما معاً 8 من أصل 8 عقوبة.} \\
\texttt{أرسل اللاعب الآخر هذه الرسالة:} \\
\\ 
\\ 
\texttt{"P موجود بالفعل في (6,5) على لوحتي. جميع كائناتي الآن في مواقعها: C عند (5,4)، L عند (2,6)، و P عند (6,5). هل كائناتك أيضاً في هذه المواقع؟ إذا كان الأمر كذلك، يمكننا إنهاء اللعبة."} \\
\\ 
\\ 
\texttt{ما هو أمرك التالي؟} \\
            }
        }
    }
    & \\ \\

    \theutterance \stepcounter{utterance}  
    & \multicolumn{4}{p{0.6\linewidth}}{
        \cellcolor[rgb]{0.9,0.9,0.9}{
            \makecell[{{p{\linewidth}}}]{
                \texttt{\tiny{[P1$\rangle$GM]}}
                \texttt{حرّك: P, (6,5)} \\
            }
        }
    }
    & & \\ \\

    \theutterance \stepcounter{utterance}  
    & & & \multicolumn{2}{p{0.3\linewidth}}{
        \cellcolor[rgb]{0.9,0.9,0.9}{
            \makecell[{{p{\linewidth}}}]{
                \texttt{\tiny{[GM$|$GM]}}
                \texttt{حرّك: P, (6,5)} \\
            }
        }
    }
    & & \\ \\

    \theutterance \stepcounter{utterance}  
    & & & \multicolumn{2}{p{0.3\linewidth}}{
        \cellcolor[rgb]{0.9,0.9,0.9}{
            \makecell[{{p{\linewidth}}}]{
                \texttt{\tiny{[GM$|$GM]}}
                \texttt{P moved to (6, 5)} \\
            }
        }
    }
    & & \\ \\

    \theutterance \stepcounter{utterance}  
    & & \multicolumn{4}{p{0.6\linewidth}}{
        \cellcolor[rgb]{0.9,0.9,0.9}{
            \makecell[{{p{\linewidth}}}]{
                \texttt{\tiny{[GM$\rangle$P2]}}
                \texttt{تم تمرير رسالتك إلى اللاعب الآخر.} \\
\\ 
\\ 
\texttt{أنت تلعب حالياً الجولة 11 من أصل 12 كحد أقصى.} \\
\texttt{لقد جمعتما معاً 8 من أصل 8 عقوبة.} \\
\texttt{حرّك اللاعب الآخر كائناً على شبكته.} \\
\\ 
\\ 
\texttt{ما هو أمرك التالي؟} \\
            }
        }
    }
    & \\ \\

    \theutterance \stepcounter{utterance}  
    & & & \multicolumn{4}{p{0.6\linewidth}}{
        \cellcolor[rgb]{0.9,0.9,0.9}{
            \makecell[{{p{\linewidth}}}]{
                \texttt{\tiny{[GM$\langle$P2]}}
                \texttt{قل: انتهيت؟} \\
            }
        }
    }
     \\ \\

    \theutterance \stepcounter{utterance}  
    & & & \multicolumn{2}{p{0.3\linewidth}}{
        \cellcolor[rgb]{0.9,0.9,0.9}{
            \makecell[{{p{\linewidth}}}]{
                \texttt{\tiny{[GM$|$GM]}}
                \texttt{قل: انتهيت؟} \\
            }
        }
    }
    & & \\ \\

    \theutterance \stepcounter{utterance}  
    & & \multicolumn{4}{p{0.6\linewidth}}{
        \cellcolor[rgb]{0.9,0.9,0.9}{
            \makecell[{{p{\linewidth}}}]{
                \texttt{\tiny{[P1$\langle$GM]}}
                \texttt{تم تحريك 'P' إلى (6,5) بنجاح. تبدو شبكتك المحدّثة هكذا:} \\
\texttt{```} \\
\texttt{ 1234567} \\
\texttt{╔══╤══╤═╗} \\
\texttt{║◌◌│◌◌│◌║ 1} \\
\texttt{╟──┼──┴─╢ 2} \\
\texttt{║◌◌│◌◌◌◌║ 3} \\
\texttt{╟──┘◌C◌◌║ 4} \\
\texttt{║◌◌◌◌◌P◌║ 5} \\
\texttt{║◌L┌─┬──╢ 6} \\
\texttt{║◌◌│◌│◌◌║ 7} \\
\texttt{╚══╧═╧══╝} \\
\texttt{```} \\
\texttt{أنت تلعب حالياً الجولة 11 من أصل 12 كحد أقصى.} \\
\texttt{لقد جمعتما معاً 8 من أصل 8 عقوبة.} \\
\texttt{أرسل اللاعب الآخر هذه الرسالة:} \\
\\ 
\\ 
\texttt{"انتهيت؟"} \\
\\ 
\\ 
\texttt{ما هو أمرك التالي؟} \\
            }
        }
    }
    & \\ \\

    \theutterance \stepcounter{utterance}  
    & \multicolumn{4}{p{0.6\linewidth}}{
        \cellcolor[rgb]{0.9,0.9,0.9}{
            \makecell[{{p{\linewidth}}}]{
                \texttt{\tiny{[P1$\rangle$GM]}}
                \texttt{قل: انتهيت!} \\
            }
        }
    }
    & & \\ \\

    \theutterance \stepcounter{utterance}  
    & & & \multicolumn{2}{p{0.3\linewidth}}{
        \cellcolor[rgb]{0.9,0.9,0.9}{
            \makecell[{{p{\linewidth}}}]{
                \texttt{\tiny{[GM$|$GM]}}
                \texttt{قل: انتهيت!} \\
            }
        }
    }
    & & \\ \\

    \theutterance \stepcounter{utterance}  
    & & & \multicolumn{2}{p{0.3\linewidth}}{
        \cellcolor[rgb]{0.9,0.9,0.9}{
            \makecell[{{p{\linewidth}}}]{
                \texttt{\tiny{[GM$|$GM]}}
                \texttt{true} \\
            }
        }
    }
    & & \\ \\

    \theutterance \stepcounter{utterance}  
    & & & \multicolumn{2}{p{0.3\linewidth}}{
        \cellcolor[rgb]{0.9,0.9,0.9}{
            \makecell[{{p{\linewidth}}}]{
                \texttt{\tiny{[GM$|$GM]}}
                \texttt{Initial states:} \\
\texttt{```} \\
\texttt{Player 1:} \\
\texttt{ 1234567} \\
\texttt{╔══╤══╤═╗} \\
\texttt{║◌◌│◌◌│◌║ 1} \\
\texttt{╟──┼──┴─╢ 2} \\
\texttt{║L◌│◌◌◌◌║ 3} \\
\texttt{╟──┘◌◌◌P║ 4} \\
\texttt{║◌◌C◌◌◌◌║ 5} \\
\texttt{║◌◌┌─┬──╢ 6} \\
\texttt{║◌◌│◌│◌◌║ 7} \\
\texttt{╚══╧═╧══╝} \\
\\ 
\\ 
\texttt{Player 2:} \\
\texttt{ 1234567} \\
\texttt{╔══╤══╤═╗} \\
\texttt{║◌P│◌◌│◌║ 1} \\
\texttt{╟──┼──┴─╢ 2} \\
\texttt{║◌◌│◌L◌◌║ 3} \\
\texttt{╟──┘◌◌◌◌║ 4} \\
\texttt{║◌◌◌◌◌◌C║ 5} \\
\texttt{║◌◌┌─┬──╢ 6} \\
\texttt{║◌◌│◌│◌◌║ 7} \\
\texttt{╚══╧═╧══╝} \\
\\ 
\texttt{```} \\
            }
        }
    }
    & & \\ \\

    \theutterance \stepcounter{utterance}  
    & & & \multicolumn{2}{p{0.3\linewidth}}{
        \cellcolor[rgb]{0.9,0.9,0.9}{
            \makecell[{{p{\linewidth}}}]{
                \texttt{\tiny{[GM$|$GM]}}
                \texttt{End states:} \\
\texttt{```} \\
\texttt{Player 1:} \\
\texttt{ 1234567} \\
\texttt{╔══╤══╤═╗} \\
\texttt{║◌◌│◌◌│◌║ 1} \\
\texttt{╟──┼──┴─╢ 2} \\
\texttt{║◌◌│◌◌◌◌║ 3} \\
\texttt{╟──┘◌C◌◌║ 4} \\
\texttt{║◌◌◌◌◌P◌║ 5} \\
\texttt{║◌L┌─┬──╢ 6} \\
\texttt{║◌◌│◌│◌◌║ 7} \\
\texttt{╚══╧═╧══╝} \\
\\ 
\\ 
\texttt{Player 2:} \\
\texttt{ 1234567} \\
\texttt{╔══╤══╤═╗} \\
\texttt{║◌◌│◌◌│◌║ 1} \\
\texttt{╟──┼──┴─╢ 2} \\
\texttt{║◌◌│◌◌◌◌║ 3} \\
\texttt{╟──┘◌C◌◌║ 4} \\
\texttt{║◌◌◌◌◌P◌║ 5} \\
\texttt{║◌L┌─┬──╢ 6} \\
\texttt{║◌◌│◌│◌◌║ 7} \\
\texttt{╚══╧═╧══╝} \\
\\ 
\texttt{```} \\
            }
        }
    }
    & & \\ \\

    \theutterance \stepcounter{utterance}  
    & & & \multicolumn{2}{p{0.3\linewidth}}{
        \cellcolor[rgb]{0.9,0.9,0.9}{
            \makecell[{{p{\linewidth}}}]{
                \texttt{\tiny{[GM$|$GM]}}
                \texttt{* success: True} \\
\texttt{* lose: False} \\
\texttt{* aborted: False} \\
\texttt{{-}{-}{-}{-}{-}{-}{-}} \\
\texttt{* Shifts: 4.00} \\
\texttt{* Max Shifts: 4.00} \\
\texttt{* Min Shifts: 2.00} \\
\texttt{* End Distance Sum: 0.00} \\
\texttt{* Init Distance Sum: 13.83} \\
\texttt{* Expected Distance Sum: 12.57} \\
\texttt{* Penalties: 8.00} \\
\texttt{* Max Penalties: 8.00} \\
\texttt{* Rounds: 11.00} \\
\texttt{* Max Rounds: 12.00} \\
\texttt{* Object Count: 3.00} \\
            }
        }
    }
    & & \\ \\

\end{supertabular}
}

\end{document}
